\documentclass[12pt, a4paper]{article}
\usepackage[utf8]{inputenc}
\usepackage[lithuanian]{babel}
\usepackage[T1]{fontenc}
\usepackage{amsmath, amssymb}
\usepackage{geometry}
\usepackage{enumitem}

\geometry{top=2cm, bottom=2cm, left=2cm, right=2cm}

\title{\textbf{Namų darbas nr. 1}}
\date{}
\begin{document}

\maketitle

\section{Sąlygos}

\begin{enumerate}[label=\textbf{\arabic* užduotis.}, leftmargin=*]
    \item Išspręskite lygtis:
    \begin{enumerate}[label=\alph*)]
        \item $(3x-1)^3-9x(1-3x)(1-x)=(x+1)^2(x+1)-x ^{3}+7 $
        \item $3x^4=2x^2+1$
        \item $3x-4\sqrt{x}=4$
        \item $\left(x^2+2x+3\right)^2=5x^2+10x+21$
    \end{enumerate}

    \item Įrodykite, jog reiškinys
    $4(x-2)^3-(2x-1)(2x+3)x+26\left(x^2-1\frac{12}{13}x\right)+33$
    gali turėti tik didžiausiąją reikšmę, o mažiausiosios neturi. Raskite tą didžiausiąją reikšmę.

    \item Duotas dvyliktojo laipsnio vienanaris
    $3x^{2m-1}\cdot y^{m(m-3)}\cdot z^{5-m}$. Nustatykite $m$ reikšmę.

    \item Duotas daugianaris $P(x)=x^2-3ax+2a+3$. Su kuria natūraliąja $a$ reikšme $P(a+2)=P(1-a)$?
\end{enumerate}

\newpage

\section{Atsakymai}

\begin{enumerate}[label=\textbf{\arabic* užduotis.}, leftmargin=*]
    \item
    \begin{enumerate}[label=\alph*)]
        \item $x=-1$ arba $x=\dfrac{3}{2}$
        \item $x=-1$ arba $x=1$
        \item $x=4$
        \item $x=-3$ arba $x=1$
    \end{enumerate}
    \item  Didžiausioji reikšmė yra $\dfrac98$.
    \item $m=4$
    \item $a=1$.
\end{enumerate}

\end{document}